% Chapter 5: Stratification and Post-Stratification in Randomized Experiments
% A First Course in Causal Inference - Peng Ding
% Rewritten in Stein motivation-first style with textbook formulas

\chapter{分层与后分层随机实验}

\section{引言：随机化的 art}

> \textit{``Block what you can and randomize what you cannot.''}
> --- Box et al. (1978, page 103)

这是 George Box 继``所有模型都是错的，但有些有用''之后第二著名的名言。本章将解释它的含义。

在完全随机实验（CRE）中，治疗分配完全由 chance 决定。这种设计在样本量趋于无穷时表现优异，但在有限样本中，治疗组和对照组可能在某些协变量上出现不平衡。这种 covariate imbalance（协变量不平衡）会使实验结果的解释变得困难——观察到的结果差异可能归因于治疗，也可能归因于协变量的不平衡。

\textbf{分层随机实验（Stratified Randomized Experiment, SRE）}正是为了主动避免这种 covariate imbalance 而设计的。

\section{SRE 的定义与记号}\label{sec:sre-definition}

考虑一个包含 $n$ 个单元的有限总体，协变量 $X_i \in \{1, \dots, K\}$ 是离散的。定义：
\[
n_{[k]} = \#\{i: X_i = k\}, \quad \pi_{[k]} = n_{[k]} / n
\]
分别为第 $k$ 层（stratum）的单元数和比例（$k = 1, \dots, K$）。

在 CRE 中，我们固定总的 $n_1$ 和 $n_0$，但各层内的分配数量 $n_{[k]1}$ 和 $n_{[k]0}$ 是随机的。有趣的是，即使在 CRE 中，各层治疗组和对照组比例之差的期望为零：
\[
\mathbb{E}\left( \frac{n_{[k]1}}{n_1} - \frac{n_{[k]0}}{n_0} \right) = 0, \quad \text{但} \quad \frac{n_{[k]1}}{n_1} - \frac{n_{[k]0}}{n_0} \neq 0 \text{ 以高概率成立}.
\label{eq:balance-discrete-CRE}
\tag{5.2}
\]

\userannotation{为什么期望为零但实际不为零？}{这正是 chance 造成的问题——虽然平均而言治疗组和对照组的协变量分布相同，但单次随机化实现中可能出现显著差异。}

\begin{Definition}[分层随机实验 (SRE)]\label{def:sre}
固定 $n_{[k]1}$ 或 $n_{[k]0}$（$k = 1, \dots, K$）。我们在离散的协变量 $X$ 的每一层内独立进行 CRE。
\end{Definition}

在农业实验中，SRE 被称为\textbf{随机区组设计（randomized block design）}，各层称为区组（blocks）。类似地，分层随机化也称为 block randomization。

\textbf{层别倾向得分（stratum propensity score）}：
\[
e_{[k]} = \frac{n_{[k]1}}{n_{[k]}}
\]
这个概念在第 \cref{sec:propensity-score} 章（第 \ref{ch:observational} 部分）讨论观测性研究时将扮演核心角色。

\textbf{SRE 与 CRE 的关键区别}：
\begin{enumerate}
  \item SRE 的所有可行随机化集合是 CRE 的可行随机化集合的子集；
  \item 在 SRE 中 $e_{[k]}$ 是固定的，但在 CRE 中是随机的。
\end{enumerate}

\section{层别平均因果效应}

对于单元 $i$，我们有潜在结果 $Y_i(1)$ 和 $Y_i(0)$，个体因果效应 $\tau_i = Y_i(1) - Y_i(0)$。第 $k$ 层的层别平均因果效应为：
\[
\tau_{[k]} = n_{[k]}^{-1} \sum_{X_i = k} \tau_i.
\]

总体平均因果效应是各层别平均因果效应的加权平均：
\[
\tau = n^{-1} \sum_{i=1}^n \tau_i = \sum_{k=1}^K \pi_{[k]} \tau_{[k]}.
\tag{5.6}
\]

如果我们关心 $\tau_{[k]}$，可以在第 $k$ 层内使用 CRE 的方法进行推断。下面讨论对 $\tau$ 的统计推断。

\section{层别化实验中的 FRT}\label{sec:sre-frt}

与 CRE 类似，我们首先讨论 SRE 中的 Fisher 随机化检验（FRT）。Sharp null 假设仍然是：
\[
H_{0\Fisher}: Y_i(1) = Y_i(0) \text{ 对所有单元 } i = 1, \dots, n.
\]

FRT 的基本思想适用于任何随机实验：我们可以用任何检验统计量 $T = T(\mathbf{Z}, \mathbf{Y}, \mathbf{X})$，其中 $\mathbf{Z}, \mathbf{Y}, \mathbf{X}$ 分别是治疗向量、观测结果向量和协变量矩阵。在 SRE 和 $H_{0\Fisher}$ 下，$T$ 具有已知分布，因为 $\mathbf{Z}$ 在 SRE 下有已知分布。

\textbf{两个关键细节}：
\begin{enumerate}
  \item 当模拟治疗向量时，必须根据 \cref{def:sre} 在 $X$ 的各层内对治疗指标进行置换（这有时被称为\textbf{条件随机化检验}或\textbf{条件置换检验}）；
  \item 应选择能够反映 SRE 本质的检验统计量。
\end{enumerate}

\subsection{检验统计量的规范选择}

\mparafh{SRE 中的差值均值估计量}

受估计 $\tau$ 的启发（见 \cref{sec:sre-neymanian}），我们可以在 FRT 中使用以下分层估计量：
\[
\hat\tau_\SRE = \sum_{k=1}^K \pi_{[k]} \hat\tau_{[k]},
\]
其中
\[
\hat\tau_{[k]} = n_{[k]1}^{-1} \sum_{i=1}^n \mathbb{I}(X_i = k, Z_i = 1) Y_i - n_{[k]0}^{-1} \sum_{i=1}^n \mathbb{I}(X_i = k, Z_i = 0) Y_i
\]
是第 $k$ 层内治疗组和对照组结果均值之差。

\mparafh{学生化分层估计量}

受简单两样本问题中学生化统计量的启发，我们可以在 FRT 中使用以下学生化统计量：
\[
t_\SRE = \frac{\hat\tau_\SRE}{\sqrt{\hat V_\SRE}},
\]
其中
\[
\hat V_\SRE = \sum_{k=1}^K \pi_{[k]}^2 \left( \frac{\hat S_{[k]}^2(1)}{n_{[k]1}} + \frac{\hat S_{[k]}^2(0)}{n_{[k]0}} \right).
\tag{5.7}
\]

这里 $\hat S_{[k]}^2(1)$ 和 $\hat S_{[k]}^2(0)$ 分别是治疗组和对照组结果在第 $k$ 层内的样本方差。这个统计量的精确形式来自下面 \cref{sec:sre-neymanian} 讨论的 Neymanian 视角。

\mparafh{合并 Wilcoxon 秩和统计量}

我们首先在第 $k$ 层内计算 Wilcoxon 秩和统计量 $W_{[k]}$（回忆 \cref{ex:wilcoxon-rank-sum}），然后合并为：
\[
W_\SRE = \sum_{k=1}^K c_{[k]} W_{[k]}.
\]

基于不同的渐近方案和最优性准则，\citet{VanElteren:1960} 提出了两种加权方法。一种使用 $c_{[k]} = 1 / (n_{[k]1} n_{[k]0})$，另一种使用 $c_{[k]} = 1 / (n_{[k]} + 1)$。

\mparafh{Hodges–Lehmann 对齐秩统计量}

Van Elteren (1960) 的统计量在层数较少且各层样本量较大时表现良好。但当层数很多且各层样本量较小时，它不能进行足够多的比较，可能丢失数据中的信息。\citet{HodgesLehmann:1962} 提出了一种在对齐后进行更多跨层比较的检验统计量。他们首先将结果中心化为：
\[
\tilde Y_i = Y_i - \bar Y_{[k]}, \quad \text{若 } X_i = k,
\]
其中 $\bar Y_{[k]} = n_{[k]}^{-1} \sum_{X_i = k} Y_i$ 是第 $k$ 层内的层别均值。然后获得池化结果 $(\tilde Y_1, \dots, \tilde Y_n)$ 的秩 $(\tilde R_1, \dots, \tilde R_n)$，最后构建检验统计量：
\[
\tilde W = \sum_{i=1}^n Z_i \tilde R_i.
\]

\mparafh{Kolmogorov–Smirnov 统计量}

我们计算 $D_{[k]}$，即第 $k$ 层内治疗组和对照组结果的经验分布函数之间的最大差距。最终的检验统计量可以是：
\[
D_\SRE = \sum_{k=1}^K c_{[k]} D_{[k]}, \quad \text{或} \quad D_{\max} = \max_{1 \leq k \leq K} c_{[k]} D_{[k]},
\]
其中 $c_{[k]} = \sqrt{n_{[k]1} n_{[k]0} / n_{[k]}}$。

\section{Neymanian 推断}\label{sec:sre-neymanian}

\subsection{点估计与区间估计}

SRE 的统计推断建立在一个事实之上：它本质上由 $K$ 个独立的 CRE 组成。基于此，我们可以将 Neyman (1923) 的结果推广到 SRE。

在第 $k$ 层内，差值均值 $\hat\tau_{[k]}$ 是 $\tau_{[k]}$ 的无偏估计，其方差为：
\[
\text{var}(\hat\tau_{[k]}) = \frac{S_{[k]}^2(1)}{n_{[k]1}} + \frac{S_{[k]}^2(0)}{n_{[k]0}} - \frac{S_{[k]}^2(\tau)}{n_{[k]}},
\tag{5.8}
\]
其中 $S_{[k]}^2(1), S_{[k]}^2(0), S_{[k]}^2(\tau)$ 分别是潜在结果和个体因果效应在第 $k$ 层内的方差。

因此，分层估计量 $\hat\tau_\SRE = \sum_{k=1}^K \pi_{[k]} \hat\tau_{[k]}$ 是 $\tau$ 的无偏估计，方差为：
\[
\text{var}(\hat\tau_\SRE) = \sum_{k=1}^K \pi_{[k]}^2 \text{var}(\hat\tau_{[k]}).
\tag{5.9}
\]

如果 $n_{[k]1} \geq 2$ 且 $n_{[k]0} \geq 2$，我们可以获得第 $k$ 层内治疗组和对照组结果的样本方差 $\hat S_{[k]}^2(1)$ 和 $\hat S_{[k]}^2(0)$，并构建一个\textbf{保守方差估计量}：
\[
\hat V_\SRE = \sum_{k=1}^K \pi_{[k]}^2 \left( \frac{\hat S_{[k]}^2(1)}{n_{[k]1}} + \frac{\hat S_{[k]}^2(0)}{n_{[k]0}} \right).
\tag{5.10}
\]

基于 $\hat\tau_\SRE$ 的正态近似，我们可以构建 $\tau$ 的 Wald 型 $1 - \alpha$ 置信区间：
\[
\hat\tau_\SRE \pm z_{1-\alpha/2} \sqrt{\hat V_\SRE}.
\tag{5.11}
\]

从假设检验的角度，在 $H_{0\Neyman}: \tau = 0$ 下，我们可以将 $t_\SRE = \hat\tau_\SRE / \sqrt{\hat V_\SRE}$ 与标准正态分位数进行比较，得到渐近 p 值。\footnote{推导见附录 \cref{sec:appendix-sre-derivation}}

\userannotation{}{$t_\SRE$ 这个统计量出现在 \cref{sec:sre-frt} 的 FRT 示例中。第 \ref{ch:unification} 章将展示在 FRT 中使用 $t_\SRE$ 可在 $H_{0\Fisher}$ 下得到有限样本精确 p 值，在 $H_{0\Neyman}$ 下得到渐近有效 p 值。}

\section{SRE 与 CRE 的比较}\label{sec:compare-sre-cre}

与 CRE 相比，SRE 有什么好处？我们从 covariate balance 的角度引入了 SRE。此外，更好的 covariate balance 会带来更好的平均因果效应估计精度。

为了进行公平比较，我们假设 $e_{[k]} = e$（对所有 $k$），这保证了：
\[
\hat\tau = \hat\tau_\SRE.
\label{eq:stratified-equal}
\tag{5.3}
\]

我们比较抽样方差。方差分析技术启发了将总方差分解为层内方差和层间方差之和：
\[
\begin{aligned}
S^2(1) & = (n-1)^{-1} \sum_{i=1}^n \{Y_i(1) - \bar Y(1)\}^2 \\
& = \sum_{k=1}^K \left[ \frac{n_{[k]}-1}{n-1} S_{[k]}^2(1) + \frac{n_{[k]}}{n-1} \{\bar Y_{[k]}(1) - \bar Y(1)\}^2 \right],
\end{aligned}
\]
类似地，
\[
\begin{aligned}
S^2(0) & = \sum_{k=1}^K \left[ \frac{n_{[k]}-1}{n-1} S_{[k]}^2(0) + \frac{n_{[k]}}{n-1} \{\bar Y_{[k]}(0) - \bar Y(0)\}^2 \right], \\
S^2(\tau) & = \sum_{k=1}^K \left[ \frac{n_{[k]}-1}{n-1} S_{[k]}^2(\tau) + \frac{n_{[k]}}{n-1} \{\tau_{[k]} - \tau\}^2 \right].
\end{aligned}
\]

CRE 下差值均值估计量的方差分解为：
\[
\begin{aligned}
\var_{\CRE}(\hat\tau) & = \frac{S^2(1)}{n_1} + \frac{S^2(0)}{n_0} - \frac{S^2(\tau)}{n} \\
& = \underbrace{\sum_{k=1}^K \left[ \frac{\pi_{[k]}}{n_1} S_{[k]}^2(1) + \frac{\pi_{[k]}}{n_0} S_{[k]}^2(0) - \frac{\pi_{[k]}}{n} S_{[k]}^2(\tau) \right]}_{\text{与 SRE 相同}} \\
& \quad + \underbrace{\sum_{k=1}^K \left[ \frac{\pi_{[k]}}{n_1} \{\bar Y_{[k]}(1) - \bar Y(1)\}^2 + \frac{\pi_{[k]}}{n_0} \{\bar Y_{[k]}(0) - \bar Y(0)\}^2 - \frac{\pi_{[k]}}{n} (\tau_{[k]} - \tau)^2 \right]}_{\text{SRE 消除的项}} .
\end{aligned}
\tag{5.12}
\]

当层样本量较大时，$\var_{\CRE}(\hat\tau)$ 与 $\var_{\SRE}(\hat\tau_\SRE)$ 之差约为：
\[
\sum_{k=1}^K \frac{\pi_{[k]}}{n} \left\{ \sqrt{\frac{n_0}{n_1}} \{\bar Y_{[k]}(1) - \bar Y(1)\} + \sqrt{\frac{n_1}{n_0}} \{\bar Y_{[k]}(0) - \bar Y(0)\} \right\}^2 \geq 0.
\label{eq:difference}
\tag{5.4}
\]

这个差值是非负的，当且仅当
\[
\sqrt{\frac{n_0}{n_1}} \{\bar Y_{[k]}(1) - \bar Y(1)\} + \sqrt{\frac{n_1}{n_0}} \{\bar Y_{[k]}(0) - \bar Y(0)\} = 0
\]
对 $k = 1, \dots, K$ 都成立时才为零。当协变量能预测潜在结果时，上述量通常不全为零，这保证了 SRE 相对于 CRE 的效率增益。\footnote{推导见附录 \cref{sec:appendix-cre-sre-comparison}}

\userannotation{为什么分层能提高效率？}{因为分层消除了层间变异（between-stratum variability）。当协变量能预测结果时，这种消除效应显著，从而减少估计量的方差。}

\section{CRE 中的后分层}\label{sec:post-stratification}

在具有离散协变量 $X$ 的 CRE 中，第 $k$ 层内接受治疗和对照的单元数是随机的。在 SRE 中，这些数字是固定的。但如果我们在给定 $\mathbf{n} = \{n_{[k]1}, n_{[k]0}\}_{k=1}^K$ 的条件下进行条件推断，那么 CRE 就变成了 SRE。

数学上，如果 $\mathbf{n}$ 的所有分量都不为零，那么：
\[
\Pr_{\CRE}(\mathbf{Z} = \mathbf{z} \mid \mathbf{n}) = \frac{\Pr_{\CRE}(\mathbf{Z} = \mathbf{z}, \mathbf{n})}{\Pr_{\CRE}(\mathbf{n})} = \frac{1}{\prod_{k=1}^K \binom{n_{[k]}}{n_{[k]1}}}.
\label{eq:condition-to-SRE}
\tag{5.5}
\]
即，$\mathbf{Z}$ 在给定 $\mathbf{n}$ 条件下从 CRE 得到的条件分布与从 SRE 得到的分布相同。

因此，在给定 $\mathbf{n}$ 的条件下，我们可以像分析 SRE 一样分析带有离散协变量 $X$ 的 CRE。FRT 变为\textbf{条件 FRT}，Neymanian 分析变为\textbf{后分层（post-stratification）}：
\[
\hat\tau_\PS = \sum_{k=1}^K \pi_{[k]} \hat\tau_{[k]},
\tag{5.13}
\]
其形式与 $\hat\tau_\SRE$ 相同。

\textbf{分层在设计阶段使用 $X$，后分层在分析阶段使用 $X$。它们是使用 $X$ 的对偶方法。}

\userannotation{SRE vs Post-stratification}{两者在渐近意义上差异很小（Miratrix et al., 2013）。但条件 FRT 通常比无条件 FRT 更有效（Hennessy et al., 2016）。}

\section{实践中的问题}\label{sec:sre-practice}

如何选择 $X$ 来构建 SRE？理论上，$X$ 应该能预测潜在结果。在某些情况下，实验者根据先导研究等对预测协变量有足够的背景知识，那么 $X$ 的选择是 straightforward 的。在其他情况下，这种背景知识不够清晰，实验者反而根据方便性选择 $X$——例如，$X$ 可以是研究区域的指示变量或学生年级。

\textbf{$K$ 的选择是一个相关问题}。理论上，如果所有层都足够大，增加 $K$ 会提高估计效率。然而，极端大的 $K$ 甚至可能降低估计效率。在模拟研究中，我们观察到增加 $K$ 的边际收益递减。经验上，$K = 5$ 通常足以获得效率增益（这个神奇的数字 5 将在第 \ref{ch:propensity-score} 章再次出现）。

\textbf{多维连续协变量怎么办？} 如果我们有先导研究，可以建立给定这些协变量时潜在结果 $Y(0)$ 的模型，然后选择 $X$ 为预测量 $\hat Y(0)$ 的离散化版本。一般地，如果我们没有这样的先导研究或者不想做 ad hoc 的离散化，我们可以使用一种更通用的策略，称为\textbf{重随机化（rerandomization）}，这将是第 \ref{ch:rerandomization} 章的主题。

\section{本章小结}\label{sec:chapter5-summary}

\begin{enumerate}
  \item \textbf{SRE 的动机}：CRE 可能产生不期望的治疗分配，通过分层可以主动避免 covariate imbalance。
  \item \textbf{SRE 定义}：在离散协变量的各层内独立进行 CRE，固定各层的治疗分配数量。
  \item \textbf{效率增益}：当协变量能预测潜在结果时，SRE 比 CRE 更有效，因为分层消除了层间变异。
  \item \textbf{FRT 和 Neymanian 推断}：都可以推广到 SRE，分别在 \cref{sec:sre-frt} 和 \cref{sec:sre-neymanian} 中讨论。
  \item \textbf{后分层}：CRE 中给定分层数量条件后，分析方法与 SRE 相同——这是使用协变量的对偶方法。
\end{enumerate}

\userannotation{Box 名言的深意}{``Block what you can and randomize what you cannot'' 暗示：对于能控制的部分（协变量）进行分层设计，对于无法控制的随机变异则依赖随机化。这正是 SRE 的核心思想。}

\section{习题}\label{sec:chapter5-exercises}

\begin{Exercise}{\ref{exr:5-1} Covariate balance in the CRE}\label{exr:5-1}
证明 \eqref{eq:balance-discrete-CRE}：在 CRE 下，
\[
\mathbb{E}\left( \frac{n_{[k]1}}{n_1} - \frac{n_{[k]0}}{n_0} \right) = 0.
\]
\end{Exercise}

\begin{Exercise}{\ref{exr:5-2} Consequence of the constant propensity score}\label{exr:5-2}
证明 \eqref{eq:stratified-equal}：在常数倾向得分 $e_{[k]} = e$（对所有 $k$）下，
\[
\hat\tau = \hat\tau_\SRE.
\]
\end{Exercise}

\begin{Exercise}{\ref{exr:5-3} Consequence of constant individual causal effects}\label{exr:5-3}
假设个体因果效应是常数 $\tau_i = \tau$（对所有 $i = 1, \ldots, n$）。考虑以下 $\tau$ 的加权估计量类：
\[
\hat\tau_w = \sum_{k=1}^K w_{[k]} \hat\tau_{[k]},
\]
其中权重 $w_{[k]}$ 对所有 $k$ 非负。

\begin{enumerate}
  \item 找出使 $\hat\tau_w$ 对 $\tau$ 无偏的 $w_{[k]}$ 条件。
  \item 在所有无偏估计量中，找出使 $\hat\tau_w$ 方差最小的权重。
\end{enumerate}
\end{Exercise}

\begin{Exercise}{\ref{exr:5-4} Compare the CRE and SRE}\label{exr:5-4}
证明 \eqref{eq:difference}：
\[
\sum_{k=1}^K \frac{\pi_{[k]}}{n} \left\{ \sqrt{\frac{n_0}{n_1}} \{\bar Y_{[k]}(1) - \bar Y(1)\} + \sqrt{\frac{n_1}{n_0}} \{\bar Y_{[k]}(0) - \bar Y(0)\} \right\}^2 \geq 0.
\]
\end{Exercise}

\begin{Exercise}{\ref{exr:5-5} From the CRE to the SRE}\label{exr:5-5}
证明 \eqref{eq:condition-to-SRE}：
\[
\Pr_{\CRE}(\mathbf{Z} = \mathbf{z} \mid \mathbf{n}) = \frac{1}{\prod_{k=1}^K \binom{n_{[k]}}{n_{[k]1}}}.
\]
\end{Exercise}

\begin{Exercise}{\ref{exr:5-6} More FRTs for Section 5.2.2}\label{exr:5-6}
使用本章讨论的其他检验统计量扩展第 5.2.2 节的分析。
\end{Exercise}

\begin{Exercise}{\ref{exr:5-7} FRT for an SRE in Imbens and Rubin (2015)}\label{exr:5-7}
Imbens and Rubin (2015) 讨论了 1985-1986 年在田纳西州进行的 Student/Teacher Achievement Ratio (Project STAR) 实验中的一个 SRE。各层对应学校，分析单位是教师或班级。治疗等于 1 表示小班（每名教师 13-17 名学生），等于 0 表示常规班级（22-25 名学生）。结果是标准化平均数学成绩。

使用 FRT 重新分析 Project STAR 数据。使用 $\hat\tau_\SRE$、$W_\SRE$ 和 $\tilde W$ 进行 FRT，并比较 p 值。

\medskip
\noindent\textbf{注：} 本书使用 $Z$ 表示治疗指标，但 Imbens and Rubin (2015) 使用 $W$。
\end{Exercise}

\begin{Exercise}{\ref{exr:5-8} A multi-center trial}\label{exr:5-8}
Gould (1998, Table 1) 报告了一个 29 个中心的多中心试验数据。这是一个以中心为层的 SRE。该试验旨在研究非那雄胺（治疗良性前列腺增生的药物）的疗效和耐受性。在 29 个中心内，患者被随机分为三组：对照组、非那雄胺 1mg 组和非那雄胺 5mg 组。

数据集提供了结局（基线到总症状评分的变化）的汇总统计。由于未报告个体水平的结果，我们无法实施 FRT。然而，Neymanian 推断只需要汇总统计量。

分别报告比较"非那雄胺 1mg"和"非那雄胺 5mg"与"对照"的点估计量和方差估计量。
\end{Exercise}

\begin{Exercise}{\ref{exr:5-9} Data re-analyses}\label{exr:5-9}
重新分析第 4 章使用的 LaLonde 数据。分别进行 Fisherian 和 Neymanian 推断。

原始实验是一个 CRE。现在我们将原始实验视为一个 SRE。

\begin{enumerate}
  \item 将实验视为按种族（黑人、西班牙裔或其他）分层的 SRE，重新分析数据。
  \item 将实验视为按婚姻状况分层的 SRE，重新分析数据。
  \item 将实验视为按高中文凭指标分层的 SRE，重新分析数据。
\end{enumerate}

与 CRE 下的结果进行比较。
\end{Exercise}
